\documentclass{article}
\usepackage{amsmath}

\begin{document}

\section*{SEA3004F -- Module 1 Examples 2}

\begin{enumerate}
    \item Write the total derivative for the density of seawater.
    
    \item Write down the expression for the coefficients in (1) and explain their impact on the density of seawater.
    
    \item The atmospheric temperature at the tropics every 500 m from sea level up to 2 km is:
    \[
    288.15,\; 284.90,\; 281.65,\; 278.40,\; 275.15 \ \text{K}
    \]
    Compute the numerical gradient of the saturation vapour pressure using the simplified equation
    \[
    e_s = A e^{\beta T}
    \]
    with
    \[
    A = 6.11 \ \text{hPa}, \quad \beta = 0.067 \ ^\circ\text{C}^{-1}.
    \]
    
    \item Which seawater temperature is indicated with the symbol $\theta$?
    
    \item Describe the main differences between atmospheric and oceanic convection processes.
    
    \item Name the forces governing the motion of sea ice.
    
    \item A planet in the Far-Far-Away galaxy, which has the mass of the Earth but a different atmospheric gas composition, has an isothermal atmosphere with temperature $-40^\circ$C. In this atmosphere, pressure decays exponentially with scale height
    \[
    H = \frac{RT}{g}.
    \]
    A radiosonde measured the pressure at 20 km height from the surface to be 100 hPa. Using a gas constant
    \[
    R = 287 \ \text{J kg}^{-1} \text{K}^{-1},
    \]
    calculate the surface pressure in hPa. Is this greater or smaller than the surface pressure of the Earth?
    
\end{enumerate}

\end{document}